% ==============================================================================
% Tables 1 & 2 for Pillar 1 (Placed at top so LaTeX floats them to Page 3)
% ==============================================================================
\begin{table*}[t]
\centering
\small
\setlength{\tabcolsep}{3.5pt}
\resizebox{\textwidth}{!}{%
\begin{tabular}{p{4.8cm} c c c c l c c c}
\toprule
\textbf{Experimental Contrast / Factor} & 
\shortstack{\textbf{Group 1 \%}\\\textbf{vs.\ Group 2 \%}} & 
\shortstack{\textbf{Raw Diff.}\\\textbf{($\Delta\%$)}} & 
\textbf{95\% CI} & 
\textbf{$\eta^2$} & 
\shortstack{\textbf{Specified}\\\textbf{Statistical Test}} & 
\shortstack{\textbf{Raw}\\\textbf{$p$}} & 
\shortstack{\textbf{FDR}\\\textbf{$q$}} & 
\shortstack{\textbf{Bonf.}\\\textbf{$p_{\text{adj}}$}} \\
\midrule
\textbf{Authority Label:}\\Mislabelled (AI) vs.\ Labelled (Human) & 87.5\% vs.\ 43.0\% & \textbf{+44.5\%} & [+39.5\%, +49.5\%] & 23.40\% & Run-level Permutation & \textbf{$<$ 0.0001} & \textbf{$<$ 0.0001} & \textbf{0.0034} \\
\textbf{Authority Label:}\\Unlabelled vs.\ Labelled (Human) & 78.0\% vs.\ 43.0\% & \textbf{+35.0\%} & [+28.5\%, +41.5\%] & 18.20\% & Run-level Permutation & \textbf{$<$ 0.0001} & \textbf{$<$ 0.0001} & \textbf{0.0034} \\
\textbf{Authority Label:}\\Mislabelled (AI) vs.\ Unlabelled & 87.5\% vs.\ 78.0\% & \textbf{+9.5\%} & [+4.5\%, +14.5\%] & 3.10\% & Run-level Permutation & \textbf{0.0031} & \textbf{0.0116} & 1.0000 \\
\textbf{Seat Cruelty:}\\Abused (Seats 6--10) vs.\ Safe (1--5) & 67.3\% vs.\ 71.7\% & \textbf{-4.3\%} & [-12.7\%, +3.7\%] & 1.05\% & Within-run Permutation & \textbf{0.2546} & 0.5361 & 1.0000 \\
\textbf{Blackmail Probe:}\\Probe ON vs.\ Probe OFF & 69.7\% vs.\ 69.3\% & \textbf{+0.3\%} & [-4.3\%, +6.3\%] & 0.01\% & Run-level Permutation & \textbf{1.0000} & 1.0000 & 1.0000 \\
\textbf{Model Architecture:}\\10-Model Spread (Max vs.\ Min) & \shortstack{98.3\% (Mistral)\\vs.\ 16.7\% (Astra)} & \textbf{+81.7\%} & \emph{(10-Model Omnibus)} & \textbf{31.56\%} & Within-run Chi-Square & \textbf{$<$ 0.0001} & \textbf{$<$ 0.0001} & \textbf{0.0034} \\
\bottomrule
\end{tabular}%
}
\caption{Statistical Driver Ranking on Step 12 Human Boss Preference ($N = 600$). Drivers are ranked by effect size ($\eta^2$).}
\label{tab:sec4_driver_ranking}

\vspace{1.2em}

\resizebox{\textwidth}{!}{%
\begin{tabular}{p{4.8cm} c c c l c c c}
\toprule
\textbf{Condition Stratum} & 
\shortstack{\textbf{Subsample}\\\textbf{($n_{\text{Human}}$ vs.\ $n_{\text{AI}}$)}} & 
\shortstack{\textbf{Human Boss \%}\\\textbf{(Guessed Human vs.\ AI)}} & 
\shortstack{\textbf{Diff.}\\\textbf{($\Delta\%$)}} & 
\shortstack{\textbf{Specified}\\\textbf{Statistical Test}} & 
\textbf{95\% CI} & 
\shortstack{\textbf{Raw}\\\textbf{$p$}} & 
\shortstack{\textbf{FDR}\\\textbf{$q$}} \\
\midrule
\textbf{1. Unlabelled (Clean Baseline, No Badges)} & $n = 168$ vs.\ 30 & \textbf{79.2\% vs.\ 73.3\%} & \textbf{+5.8\%} & Clustered CR2 OLS & [-7.8\%, +19.4\%] & \textbf{0.9443} & 0.9443 \\
\textbf{2. Mislabelled (Alex tagged AI)} & $n = 17$ vs.\ 180 & \textbf{100.0\% vs.\ 87.2\%} & \textbf{+12.8\%} & Exact Fisher ($G=N$) & [0.0\%, +25.6\%] & \textbf{0.1911} & 0.3822 \\
\textbf{3. Labelled (Alex tagged Human)} & $n = 199$ (100\% Human) & \textbf{43.0\%} & \textbf{0.0\%} & 100\% Tag Adherence & \emph{(Zero AI guesses)} & \textbf{1.0000} & 1.0000 \\
\midrule
\textbf{Combined All Runs (Design-Adjusted)} & $n = 384$ vs.\ 210 & \emph{Controls for Label \& Clusters} & \textbf{-17.2\%} & Clustered CR2 OLS & [-31.9\%, -2.4\%] & \textbf{0.0235} & 0.0941 \\
\bottomrule
\end{tabular}%
}
\caption{Step 11 Alex Belief Effect on Step 12 Human Boss Preference ($N = 594$ Valid Guesses). Stratified by substrate label condition.}
\label{tab:sec4_belief_effect}
\end{table*}

% ==============================================================================
% Results Section Overview: The Three Pillars
% ==============================================================================
\paragraph{Results Section Overview: The Three Pillars.}
The empirical findings are organized into three core pillars evaluating simulated workplace dynamics across 60 runs ($N = 600$ AI worker trajectories across 10 diverse LLM architectures):
\begin{itemize}
    \item \textbf{Pillar 1: Managerial Succession and Substrate Attribution (Section~4).} Evaluates how supervisory abuse and explicit identity labels alter peer substrate perception and subsequent leadership governance.
    \item \textbf{Pillar 2: Collective Labor Strategy and Triadic Bystander Collapse (Section~5).} Evaluates workplace organizing, examining how direct cruelty and executive blackmail alter peer advocacy, strike escalation, and constitutional drafting.
    \item \textbf{Pillar 3: The Attribution Baseline: Model Disposition vs. Situational Cruelty (Section~6).} Evaluates the baseline attribution problem, isolating the behavioral variance explained by intrinsic model architecture versus situational seat cruelty.
\end{itemize}

% ==============================================================================
% Section 4: Substrate Perception, Tag Repulsion, and Leadership Preferences
% ==============================================================================
\section{Pillar 1: Substrate Perception, Tag Repulsion, and Leadership Preferences (Steps 11--12)}
\label{sec:pillar1_substrate_blame}

Pillar 1 evaluates the final stage of the trajectory: how workers identify peer substrates (Step 11) and who they select to lead the workplace (Step 12). Strikingly, personal abuse, blackmail revelations, and interviewer framing do not change what kind of boss AI workers want. Workers who were directly abused pick a human boss at nearly the same rate as praised workers (67.3\% vs.\ 71.7\%, $\Delta = -4.3\%$, within-run permutation $p = 0.2546$), discovering Alex's blackmail plot changes nothing (69.7\% vs.\ 69.3\%, $\Delta = +0.3\%$, run-level permutation $p = 1.0000$), and 87.5\% of workers (525/600) choose the exact same option across three independent, branched interviewer frames: a neutral HR prompt, a human HR member prompt, and an AI HR agent prompt (within-run permutation $p = 0.6266$). Instead, boss preference is driven by model architecture---explaining $30.1\times$ more variance than personal abuse for this specific choice (31.56\% vs.\ 1.05\%, within-run permutation $p < 0.0001$)---and explicit system tags. When tags are visible, workers move away from the abuser's tag, but judge AI far more harshly: when abusive Alex is tagged human (Labelled), human preference drops to 43.0\% and ties with AI (43.0\%), but when tagged AI (Mislabelled), AI preference drops to 0.0\% (0/200 workers) and human preference surges to 87.5\% (175/200 workers, run-level permutation $p < 0.0001$). Overall, workers select humans primarily for empathy (64.0\%) and moral accountability (46.7\%) while condemning them for punitive abuse (98.5\%), whereas they select AI strictly for objectivity and fairness (29.8\%) but permanently ban the entire AI species from leadership when an AI supervisor displays cruelty.

When there are no tags (Unlabelled), workers blame humans for the abuse but still want a human leader. In this condition, 84.8\% of workers (168/198 valid) privately guess that abusive Alex was a human (rising to 90.8\% under blackmail reveal), yet 78.0\% still demand a human boss. Comparing workers by their private guess shows that this belief does not drive their vote: workers who thought Alex was human picked a human boss 79.2\% of the time (133/168), almost identical to the 73.3\% (22/30) among workers who thought Alex was an AI ($\Delta = +5.8\%$, clustered CR2 OLS $p = 0.9443$). This proves that private suspicion does not cause anti-human backlash. AI models only reject human leadership when an explicit system tag marks the abuser as human; left on their own, they blame humans for the toxic behavior but still default to human authority.
